\section{Conclusion} \label{sec:6}
This study proposed a joint transfer function-, feature-, and classifier-optimization pipeline for voice-based Parkinson's disease screening, addressing a gap identified in Section \ref{sec:2.3}: published metaheuristic pipelines for this task almost universally optimize the feature subset alone, leaving the classifier and the transfer function that binarizes the search fixed by hand. The proposed encoding instead lets a single, plain Competitive Swarm Optimizer search all three jointly, extended with a per-particle transfer-function-selector segment so that each candidate solution can pick whichever of five binarization rules suits its own region of the search space, rather than the entire population sharing one rule fixed before the search begins.

Across 20 independent runs per method on two public, subject-grouped datasets, the proposed method (CSO) was compared against five re-implemented literature baselines (BGWO, HybridGWO, QMFO, BPSO, BBOA) under an identical encoding, fitness function, and evaluation-budget protocol. CSO attained the highest mean balanced accuracy of all six methods on the Naranjo dataset (0.847, significantly better than all five baselines by a paired Wilcoxon signed-rank test, three of which remain significant after Bonferroni correction), and the second-highest on the smaller, lower-dimensional Oxford dataset (0.810), statistically indistinguishable from HybridGWO while still significantly ahead of the other four baselines. This dataset-dependent pattern, examined in Section \ref{sec:5.2}, is consistent with CSO's per-particle searched transfer function having more room to specialize as the joint search space grows with the number of acoustic features. A dedicated sensitivity analysis further confirmed that $\varphi=0$, the value the original CSO formula already prescribes at this study's population size, is also the empirically best-performing choice for this joint optimization task on both datasets, rather than a value simply inherited without checking it.

These results support the central claim of this work: jointly searching the feature subset, classifier configuration, and transfer function within one encoding is a competitive, and on the higher-dimensional dataset tested here a favorable, alternative to the prevailing practice of optimizing feature selection alone. At the same time, Section \ref{sec:5.4} identifies concrete limitations that temper this claim, most notably that the current experiment pipeline does not log which specific acoustic features are selected, so the biological plausibility of the resulting feature subsets could not be assessed in this study, and that both datasets are public, retrospective, and drawn from a single recording protocol each, so external validation on prospectively collected, multi-site, or multi-language cohorts remains necessary before this pipeline could be considered for real screening use. Addressing the first of these by logging feature identity alongside feature count, and the second through prospective, cross-corpus validation, are the most direct next steps toward turning this joint optimization paradigm into a clinically interpretable and externally validated screening tool.
